\chapter{结论与建议}

\section{研究结论}

本研究以泉州世界遗产为课程资源，以项目化学习为组织方式，在初中语文综合性学习中开展了两轮行动研究。通过对问卷、课堂观察、学生作品与访谈资料的分析，可以得出以下结论。

第一，地方世界遗产可以转化为具有语文学习价值的项目资源。转化的关键不在于资源本身是否“有名”，而在于能否从资源中提炼出与学生经验相连、需要语言实践才能完成的问题。本研究以“如何让更多人读懂泉州世界遗产”为驱动性问题，将遗产信息、现场观察、人物访谈与解说词写作组织成任务链条，使学生在真实表达中综合运用阅读、写作、口语交际与探究能力。实践证明，这种转化是可行的，也是必要的。

第二，项目化学习能够提升学生参与综合性学习的主动性与文化理解水平。问卷结果显示，学生在文化理解、语言运用、合作探究与问题解决四个维度上的后测得分均显著高于前测，其中文化理解维度的提升最为明显。访谈与作品分析进一步表明，走访、访谈与公开成果让学生对遗产产生了更具体的认识，也让他们感受到语文学习与真实生活的联系。学生不再把综合性学习视为“查资料、做手抄报”的额外任务，而是把它看作一次需要合作、探究与表达的真实实践。

第三，项目化学习的实施需要精细的支架与持续的评价。第一轮实施中出现的主题过宽、资料堆砌、合作不均衡等问题，说明项目化学习并不会自动带来高质量学习。只有通过缩小主题、提供资料筛选工具、设置角色卡、将评价量表转化为学生可理解的问题清单，项目才能从“热闹的活动”走向“有深度的学习”。教师的角色不是退居一旁，而是在关键节点提供示范、反馈与支持。

第四，项目化学习在语文学科中的实施存在需要正视的张力。项目逻辑要求持续探究与真实成果，课时逻辑要求在规定时间内完成教学任务；学生自主要求开放空间，教师指导要求必要约束；过程评价要求丰富证据，结果评价要求可操作性。这些张力无法通过一次项目彻底解决，而需要教师在具体情境中不断权衡与调整。

\section{教学建议}

基于研究结论，本文提出以下建议。

对教师而言，首先应建立“目标—任务—证据”一致的设计意识。设计项目时，先明确学生应达到的语文学习目标，再思考哪些任务能够承载这些目标，最后确定哪些证据能够证明学习发生。避免先设计活动、再补目标的逆向操作。其次，应为学生提供可操作的支架。支架不是降低要求，而是帮助学生跨越当前能力与目标之间的差距。资料筛选表、访谈提纲、写作结构图、评价问题清单等工具，都可以根据学情进行调整。再次，应重视项目过程中的语言训练。项目化学习不应挤占语言文字学习的时间，而应把阅读、写作、口语交际嵌入任务之中，让学生在“做事”的同时“学语文”。

对学校而言，可以建立地方文化课程资源库，整合遗产点介绍、地方文献、口述史、影像资料与教学案例，降低教师开发资源的成本。学校还可以通过跨学科教研、馆校合作与社会实践基地建设，为项目化学习提供支持。与此同时，学校应给教师一定的课程自主空间，允许教师在课时安排、作业设计与评价方式上作出调整，避免用统一的进度与考试要求束缚项目化学习的开展。

对课程开发者而言，应警惕把地方文化资源“景点化”“标签化”的倾向。泉州世界遗产不是一组供人参观的景点，而是与社区生活、城市发展、文化传承紧密相连的活的遗产。课程开发应引导学生关注遗产的当代处境，讨论保护与利用、传承与创新之间的关系，培养学生的文化判断力与责任感。只有这样，地方文化资源才能真正进入语文课程，而不是成为一次性的活动素材。

\section{研究不足与展望}

本研究存在一些不足。第一，研究样本较小，仅在一所学校的一个班级开展，研究结论的推广需要谨慎。第二，研究周期较短，两轮行动研究共八周，难以考察项目化学习对学生文化理解与语言运用的长期影响。第三，研究未设置对照班，无法完全排除其他因素对结果的影响。第四，研究者既是教学设计者又是实施者与评价者，可能存在研究者偏见。为减少这一影响，研究采用了多种数据来源与同伴互评，但仍需在后续研究中加以改进。

未来研究可以从三个方向拓展。其一，扩大样本范围，在不同学校、不同年级开展比较研究，考察项目化学习在不同学情下的适应性与差异性。其二，延长研究周期，采用追踪研究设计，观察学生文化理解与语言能力的变化过程。其三，引入更严格的实验或准实验设计，结合标准化测评与质性分析，进一步检验项目化学习的效果。此外，还可以研究数字技术在地方遗产教育中的作用，例如利用虚拟现实、音频导览与社交媒体，拓展学生接触遗产、表达理解的方式。

总之，地方世界遗产为语文综合性学习提供了丰富而独特的课程资源，项目化学习为这些资源的转化提供了可行的组织方式。本研究是一次初步的探索，期待更多研究者与一线教师加入，共同推动地方文化资源与语文课程的深度融合。
